\documentclass{article}

\usepackage{fullpage}
\usepackage{url}
\usepackage{graphicx}
\usepackage{amsmath}
\usepackage{physics}

\title{Quantum Mechanics II Excercise 1}
\author{Idan Stark}
\date{\today}

\begin{document}
\maketitle
\section{Natrual Units}
\label{sec:ex1}
\subsection{Unit Equivilence}
In CGS, the speed of light is measured in $cm/sec$. Therefore, for a unit system where $c = 1$, we get:
\begin{equation}
    c = 3\cdot10^{10}\:cm / sec = 1
\end{equation}
\begin{equation}
    1 sec = 3\cdot10^{10}\:cm
\end{equation}
Therefore, distance and time are measured using the same units. Using the same logic, in CGS, mass is measured in grams and energy is measured in ergs. Assigning equation (2) into the definition of 1 erg, we get:
\begin{equation}
    1\:erg = 1\:g\cdot\frac{cm^2}{sec^2} = 1\:g\cdot\frac{cm^2}{9\cdot10^{20}\:cm^2} = \frac{1}{9}\cdot10^{-20}\:g
\end{equation}
And so we got that, assuming $c = 1$, the units for mass and energy are unified. This can also be shown using the relation from special relativity $E = mc^2$. As instructed, the units of mass and energy will be $eV$.
\newline
Now let's use the fact that $\hbar = 1$. $\hbar$ Has units of angualr momentum. Angular momentum has units of $energy \times time$. Since $\hbar$ is unitless, this means that angular momentum is unitless, and requires the units of energy to be opposite to the unit of time:
\begin{equation}
    \hbar = 10^{-27}\:erg\cdot sec = 1
\end{equation}
\begin{equation}
    1\:sec = 10^{27}\:erg^{-1}
\end{equation}
Therefore the units of distance and time are opposite to the unit of energy and mass, and $eV^{-1}$ can be used to measure them. Writing all of the units in ${eV}$ is now just an issue of conversion from erg. 
\begin{equation}
    1\:erg = 6.242\cdot10^{11}\:eV
\end{equation}
\begin{equation}
    1\:g = 9\cdot10^{20}\:erg = 5.617\cdot10^{32}\:eV
\end{equation}
\begin{equation}
    1\:sec = 10^{27}\:erg^{-1} = 1.6\cdot10^{15}\:eV^{-1}
\end{equation}
\begin{equation}
    1\:cm = sec/(3\cdot10^{10}) = 5.34\cdot10^{4}\:eV^{-1}
\end{equation}
\subsection{Newton's Constant and Plank Mass}
Using the values from equations (6) to (9):
\begin{equation}
    G_N = 6.7\cdot10^{-8} cm^3sec^{-2}g^{-1}
\end{equation}
\begin{equation}
    G_N = 6.7\cdot10^{-8}\cdot(5.34\cdot10^{4}\:eV^{-1})^{3}\cdot(1.6\cdot10^{15}\:eV^{-1})^{-2}\cdot(5.617\cdot10^{32}\:eV)^{-1}
\end{equation}
\begin{equation}
    G_N = 7.095\cdot10^{-57}\:eV^{-2}
\end{equation}
This is the very close to the value we would get using Plank Mass $M_{pl} = 1.2\cdot10^{19}\:GeV$:
\begin{equation}
    G_N = M_{pl}^{-2} = (1.2\cdot10^{19}\:GeV)^{-2} = 6.94\cdot10^{-57}\:eV^{-2}
\end{equation}
With a difference of around $2.1\%$, resulting probably from our rounding of the numbers.
\subsection{The Actions of different systems}
\subsubsection{Tennis ball in a court}
The average serve velocity is around $174 \ km/h = 48.33 \ m/s $\footnote{https://toomanyrackets.com/average-serve-speeds-in-tennis/}. The length of a tennis court is $23.78$ meters. Assuming the ball is being served horizontally, the time it would take to complete that distance is $0.492$ seconds. In that time, the ball drops
\begin{equation}
    \Delta y = \frac{1}{2}gt^2 = 1.18 \ m
\end{equation}
Assuming the pitch is done from an average height of 1.5 meters, we can ignore the floor. We will also ignore air resistance, even though it affects the motion, it will not make a visible change in the action relative to that of the water molecule. The speed of the ball is then, as a function of time:
\begin{equation}
    \dot{x}^2 = v_0^2 + (gt)^2 = 2335.79 + 96.04t^2
\end{equation}
\begin{equation}
\begin{gathered}
    \frac{1}{\hbar}S[x_c(t)] = \\
    = \frac{1}{\hbar} \int_{0}^{0.492 \ sec} \frac{1}{2} 0.056 \ kg \cdot \left( 2335.79 + 96.04t^2\right)dt \sim 3.2 \cdot 10^{35}
\end{gathered}
\end{equation}
\subsubsection{Water molecule in a teacup}
The mean free path of a water molecule is roughly $2.5$ \AA \footnote{https://galileo.phys.virginia.edu/classes/304/h2o.pdf}. During this time, the molecule moves in an average velocity depending on the temperature, following the formula
\begin{equation}
    v = \sqrt{\frac{2 k_B T}{m}}
\end{equation}
Assuming the temperature of the tea is roughly $60^\circ C$ and the mass of a water molecule is $3 \cdot 10^{-26} \ kg$, we get and average speed of $553.75 \ m/s$. This makes the mean free time $T = 4.5 \cdot 10^{-13}$ seconds. Assigning all of these into the action formule, we get:
\begin{equation}
\begin{gathered}
    \frac{1}{\hbar}S[x_c(t)] = \\
    = \frac{1}{\hbar} \int_{0}^{4.5\cdot10^{-13}} \frac{1}{2} 10^{-26} \ kg \cdot \left(553.75 \ m/s\right)^2 \ dt \sim 21
\end{gathered}
\end{equation}
\section{Smeared Amplitude}

\end{document}